\section{Requerimientos Específicos}

Para la priorización de los requisitos listados en las siguientes secciones, se ha adoptado formalmente la técnica \textbf{MoSCoW,} un método de priorización estándar en marcos de trabajo ágiles (DSDM Consortium, 2014). Este método clasifica los requisitos según su nivel de criticidad para el éxito del Producto Mínimo Viable (MVP). La columna \textbf{``Prioridad''} en las tablas refleja directamente esta clasificación, utilizando las siguientes etiquetas estándar:

\begin{itemize}
    \item \textbf{M (Must Have):} Requisitos no negociables, absolutamente esenciales para la operación básica del sistema.
    \item \textbf{S (Should Have):} Requisitos importantes que añaden un valor significativo y deberían incluirse si es posible.
    \item \textbf{C (Could Have):} Requisitos deseables o mejoras de calidad de vida que se incluirán solo si hay tiempo y recursos de sobra.
    \item \textbf{W (Won't Have):} Requisitos que se ha acordado explícitamente \textbf{no} incluir en esta versión, pero que se registran para futuras consideraciones.
\end{itemize}

Este enfoque asegura que el esfuerzo de desarrollo se concentre en entregar un producto funcional y de valor desde el primer sprint.

\subsection{Requisitos de Interfaz}

\subsubsection{Interfaces de Usuario}
\begin{itemize}
    \item El sistema debe proporcionar una interfaz web responsiva basada en estándares \textbf{HTML5/CSS3} para el módulo de gestión, asegurando su visualización correcta en monitores de escritorio con resolución mínima de 1366x768.
    \item El sistema debe proporcionar una interfaz táctil optimizada para dispositivos móviles (Tablets) con sistemas operativos \textbf{Android 10+}, soportando orientación vertical y horizontal.
\end{itemize}

\subsubsection{Interfaces de Hardware}
\begin{itemize}
    \item El sistema debe soportar la interacción con las \textbf{cámaras integradas} de los dispositivos móviles para la lectura de códigos QR y captura de evidencia.
    \item El sistema debe mantener una conexión de datos con \textbf{controladores IoT (microcontroladores)} para la recepción de tramas de telemetría.
\end{itemize}

\subsubsection{Interfaces de Software}
\begin{itemize}
    \item El sistema debe exponer y consumir servicios web (APIs) utilizando el estándar \textbf{REST} y el formato de intercambio de datos \textbf{JSON} para la integración con sistemas EAM/ERP externos.
    \item El sistema debe ser capaz de interpretar archivos de modelos 3D en formatos estándar de la industria (ej. \textbf{.obj, .glb, .step}) sin requerir plugins propietarios en el navegador.
\end{itemize}

\subsubsection{Interfaces de Comunicación}
\begin{itemize}
    \item Todas las comunicaciones externas deben realizarse sobre el protocolo \textbf{HTTPS (Puerto 443)} utilizando cifrado TLS 1.2 o superior.
    \item La transmisión de datos de sensores IoT debe soportar el protocolo \textbf{MQTT} para optimizar el ancho de banda en entornos industriales.
\end{itemize}

\subsection{Requisitos Funcionales (RF)}
\footnotesize
\begin{longtable}{|p{2.2cm}|p{8.5cm}|c|p{3.5cm}|}
\hline
\textbf{Código} & \textbf{Descripción} & \textbf{Prioridad} & \textbf{HU Origen} \\ \hline
\endhead
RF-MTTO-001 & El sistema debe permitir al Planificador crear y programar Órdenes de Trabajo (WO) preventivas de forma recurrente (ej. mensual, trimestral). & M & Programación de Mantenimiento Preventivo \\ \hline
RF-VIS-002 & El sistema debe superponer iconos interactivos sobre el plano que indiquen la ubicación de los Permisos de Trabajo activos. & S & Visualización de Capa de Permisos de Trabajo \\ \hline
RF-VIS-003 & El sistema debe permitir la navegación (Zoom, Paneo, Rotación) del modelo 3D mediante mouse y gestos táctiles. & M & Navegación con Zoom en Mapa de Planta \\ \hline
RF-VIS-004 & El sistema debe generar una vista explosionada interactiva del activo 3D seleccionado. & S & Visualización de Despiece (Vista Explosionada) \\ \hline
RF-VIS-005 & El sistema debe visualizar gráficamente la ruta de aislamiento (LOTO) conectando el activo con sus puntos de corte de energía. & C & Visualización de Rutas de Aislamiento (LOTO) \\ \hline
RF-VIS-006 & El sistema debe proveer una interfaz para construir la estructura jerárquica de los activos (Padre-Hijo) según la norma ISO 14224. & M & Configuración de Jerarquía de Componentes 3D \\ \hline
RF-VIS-007 & El sistema debe permitir la asociación y carga de archivos de modelos 3D (ej. .obj, .step) a nodos específicos de la jerarquía. & S & Configuración de Jerarquía de Componentes 3D \\ \hline
RF-MTTO-010 & El sistema debe validar que todos los campos obligatorios del formulario de cierre estén diligenciados antes de permitir finalizar la orden. & S & Cierre Digital de OT en Móvil \\ \hline
RF-MTTO-011 & El sistema debe mostrar una advertencia visual bloqueante en la Orden de Trabajo si los sensores indican que el equipo está energizado o en marcha. & C & Cierre Digital de OT en Móvil \\ \hline
RF-MTTO-020 & El sistema debe permitir la consulta y descarga de manuales técnicos y planos directamente desde la interfaz móvil. & M & Consulta de Manuales en Móvil \\ \hline
RF-MTTO-021 & El sistema debe permitir al Supervisor asignar un nivel de prioridad (Alta, Media, Baja) a la Orden de Trabajo antes de aprobarla o asignarla. & S & Extracción Automática de Datos de Informes \\ \hline
RF-ADM-022 & El sistema debe analizar el histórico de fallas y sugerir ajustes automáticos a la frecuencia de mantenimiento preventivo. & C & Sugerencias de Optimización (PMO) \\ \hline
RF-INV-027 & El sistema debe permitir al Ingeniero de Proyectos crear la ficha digital de un nuevo activo. & M & Creación de Ficha Técnica de Activo \\ \hline
RF-INV-028 & El sistema debe solicitar la captura de los siguientes campos obligatorios al crear la ficha de activo: Nombre, Serial, Fecha de Compra y Garantía. & S & Creación de Ficha Técnica de Activo \\ \hline
RF-INV-029 & El sistema debe permitir al Jefe de Almacén ajustar manualmente el nivel de stock de un repuesto. & M & Ajuste Manual de Stock de Repuestos \\ \hline
RF-INV-030 & El sistema debe registrar el motivo del ajuste de stock, ya sea por ``entrada por compra'' o por ``ajuste de inventario''. & S & Ajuste Manual de Stock de Repuestos \\ \hline
RF-INV-031 & El sistema debe permitir al Ingeniero de Proyectos revisar y completar la ficha técnica de los activos que se encuentren en estado ``Borrador'' o ``En Cuarentena'' y permitir el cambio de estado del activo a `Operativo' una vez validada la información. & S & Validación y Alta de Activos Nuevos \\ \hline
RF-INV-032 & El sistema debe enviar una notificación automática al Jefe de Almacén cuando el stock de un ``Repuesto Crítico'' descienda por debajo del mínimo establecido. & S & Alertas de Stock Crítico \\ \hline
RF-INV-033 & El sistema debe permitir al Jefe de Almacén configurar el nivel de stock mínimo para cada ``Repuesto Crítico''. & S & Alertas de Stock Crítico \\ \hline
RF-ADM-036 & El sistema debe validar las credenciales del usuario para permitir el acceso a la plataforma. & M & Gestión de Usuarios y Acceso \\ \hline
RF-ADM-037 & El sistema debe mostrar el dashboard con los indicadores clave de mantenimiento una vez el usuario ingrese al módulo de KPIs. & M & Dashboard Ejecutivo de KPIs \\ \hline
RF-ADM-038 & El sistema debe permitir definir, crear, editar y eliminar roles de usuario, así como asignar permisos específicos para controlar qué acciones puede realizar cada tipo de usuario. & M & Gestión de Roles y Permisos \\ \hline
RF-ADM-039 & El sistema debe permitir crear cuentas nuevas de usuario, asignarles un rol existente y administrar sus credenciales de acceso. & M & Gestión de Usuarios y Acceso \\ \hline
RF-VIS-041 & El sistema debe validar que un activo hijo herede automáticamente la ubicación geográfica de su activo padre al ser creado. & S & Configuración de Jerarquía de Componentes 3D \\ \hline
RF-MTTO-042 & El sistema debe permitir al usuario con rol `Contratista' subir archivos (PDF/Imagen) asociados a una Orden de Trabajo. & S & Carga de Informe por Contratista \\ \hline
RF-MTTO-043 & El sistema debe procesar los documentos subidos y extraer automáticamente los metadatos clave (fecha, repuestos, horas). & C & Extracción Automática de Datos de Informes \\ \hline
RF-ADM-044 & El sistema debe permitir al Administrador crear, editar y publicar mensajes de anuncio visibles para todos los usuarios. & C & Publicación de Anuncios del Sistema \\ \hline
RF-ADM-045 & El sistema debe generar un archivo PDF descargable con el resumen de los indicadores del periodo seleccionado. & S & Exportación de Reportes en PDF \\ \hline
RF-ADM-046 & El sistema debe exponer un endpoint de datos (API/CSV) que permita la conexión con herramientas externas de Business Intelligence. & C & Sincronización de Datos con Herramientas BI \\ \hline
RF-ADM-047 & El sistema debe permitir al Administrador cargar una imagen (Logo) para personalizar el encabezado de la interfaz gráfica. & C & Configuración General del Sistema \\ \hline
RF-ADM-048 & El sistema debe permitir al Administrador seleccionar la zona horaria (Timezone) global que regirá todas las fechas y horas mostradas en la plataforma. & C & Configuración General del Sistema \\ \hline
RF-ADM-049 & El sistema debe permitir la configuración de las unidades de medida predeterminadas (Métrico/Imperial) para los datos de los activos. & C & Configuración General del Sistema \\ \hline
\end{longtable}
\normalsize

\subsection{Requisitos No Funcionales (RNF)}
\footnotesize
\begin{longtable}{|p{2.5cm}|p{3.5cm}|p{7.5cm}|c|}
\hline
\textbf{Código} & \textbf{Categoría} & \textbf{Descripción} & \textbf{Prioridad} \\ \hline
\endhead
RNF-ADM-018 & Flexibilidad & La aplicación web debe ser compatible y funcional en las versiones estables actuales de Chrome, Firefox y Edge. & M \\ \hline
RNF-ADM-019 & Mantenibilidad & El código fuente del sistema debe seguir estándares de ``Clean Code'' y estar documentado para facilitar la transferencia de conocimiento. & S \\ \hline
RNF-INV-035 & Seguridad & El sistema debe registrar un historial de auditoría inmutable para cada ajuste de stock, indicando el usuario, la fecha y el valor original/nuevo del stock. & S \\ \hline
RNF-MTTO-025 & Seguridad & El sistema debe restringir la carga de archivos a formatos seguros (PDF, JPG, PNG) y limitar el tamaño a 10 MB. & S \\ \hline
RNF-ADM-015 & Seguridad & El sistema debe cifrar toda la comunicación entre cliente y servidor utilizando el protocolo TLS 1.2 o superior. & M \\ \hline
RNF-ADM-014 & Seguridad & El sistema debe invalidar el token de sesión inmediatamente después de que el usuario cierre sesión. & S \\ \hline
RNF-ADM-013 & Seguridad & El sistema debe forzar contraseñas con longitud mínima de 12 caracteres y almacenarlas usando algoritmos de hash robustos (ej. bcrypt). & M \\ \hline
RNF-MTTO-023 & Fiabilidad & El sistema debe permitir la operación en modo offline (sin conexión) del módulo móvil, almacenando los datos localmente y sincronizándolos automáticamente cuando se restablezca la red. & M \\ \hline
RNF-ADM-017 & Fiabilidad & El sistema debe garantizar una disponibilidad operativa del 99.5\% durante el horario laboral crítico (6am - 10pm). & M \\ \hline
RNF-ADM-050 & Eficiencia de Desempeño & El sistema debe validar que el archivo de imagen del logo cumpla con las restricciones de formato (PNG/JPG) y dimensiones máximas (ej. 200x50px) para no romper el diseño de la interfaz. & C \\ \hline
RNF-ADM-016 & Eficiencia de Desempeño & El sistema debe resolver las consultas de datos y cargas de interfaz en menos de 3 segundos para el 90\% de las peticiones. & S \\ \hline
RNF-VIS-008 & Eficiencia de Desempeño & El sistema debe implementar carga bajo demanda (Lazy Loading) o Nivel de Detalle (LOD) para optimizar la visualización 3D en móviles. & S \\ \hline
RNF-VIS-024 & Compatibilidad & El modelo de datos debe soportar una jerarquía recursiva de al menos 6 niveles para cumplir con la taxonomía ISO 14224. & M \\ \hline
RNF-ADM-012 & Compatibilidad & El sistema debe implementar protocolos estandarizados (MQTT, OPC UA) para la comunicación con sensores IoT. & C \\ \hline
RNF-ADM-051 & Capacidad de Interacción & Los cambios en la configuración global (Logo, Zona Horaria) deben aplicarse en tiempo real para todos los usuarios sin requerir un reinicio del servidor. & C \\ \hline
RNF-INV-034 & Capacidad de Interacción & El sistema debe permitir completar el formulario de creación de activo en una sola pantalla sin necesidad de recargar la página. & C \\ \hline
RNF-VIS-009 & Capacidad de Interacción & La interfaz móvil debe ser responsiva (responsive), adaptando tablas y gráficos complejos a pantallas pequeñas sin perder legibilidad. & M \\ \hline
\end{longtable}
\normalsize
